% DR3 replacement text for the visual_grasp robot-operation backend.
% Author/responsible team member: Runhan Wang.
%
% Intended use:
%   1. Replace the current report subsections
%      "Inverse Kinematics and Motion Planning" and
%      "MuJoCo and Piper Robotic-Arm Platform" with this file.
%   2. The report preamble already loads amsmath, amssymb, longtable,
%      booktabs, array, graphicx, and url, which are sufficient here.
%   3. Do NOT copy the host-side task-duration values into the
%      OpenHarmony-to-host communication-latency table. They measure
%      different system boundaries.
%
% Suggested replacements in the earlier Current DR3 Status table:
% Inverse kinematics and robot model & Completed and regression-tested &
% Runhan Wang: DH-based Piper FK/IK, DH-matched MuJoCo model, and
% 483/500 no-warm-start IK regression cases solved. \\
% MuJoCo visual-grasp backend & Completed for simulation; perception
% robustness characterized & Runhan Wang: structured multi-task executor,
% YOLO/RGB-D and sim_gt perception backends, motion-plan validation, and
% machine-readable benchmark outputs. \\
% Physical Piper backend & Implemented and offline-tested; hardware acceptance
% pending & Runhan Wang: feedback-guarded continuous-waypoint backend and
% double authorization gate; no CAN/SDK motion or physical-arm acceptance is
% claimed. \\
%
% Suggested owner-table replacement:
% Simulation/physical demo hardening & Runhan Wang \\

\subsection{Inverse Kinematics and Motion Planning}
\label{sec:visual-grasp-ik-motion}

The robot-operation backend is implemented in the \texttt{visual\_grasp}
repository and is responsible for converting a bounded task command into
perception, inverse kinematics (IK), a validated motion plan, and either
MuJoCo execution or an explicitly gated Piper-hardware execution request.
The backend does not own the OpenHarmony transport, ROS~1--ROS~2 topic
mapping, or network-latency measurement. Its external command boundary is a
versioned JSON schema, \texttt{visual\_grasp.bridge.v1}.

\subsubsection{Kinematic model}

The manipulator is the six-degree-of-freedom AgileX Piper with a parallel
gripper.  The implementation uses the following homogeneous transform for
joint~$i$:

\begin{equation}
{}^{i-1}\!A_i =
\begin{bmatrix}
\cos\theta_i & -\sin\theta_i & 0 & a_i \\
\cos\alpha_i\sin\theta_i & \cos\alpha_i\cos\theta_i & -\sin\alpha_i & -d_i\sin\alpha_i \\
\sin\alpha_i\sin\theta_i & \sin\alpha_i\cos\theta_i & \cos\alpha_i & d_i\cos\alpha_i \\
0 & 0 & 0 & 1
\end{bmatrix},
\qquad
\theta_i=q_i+\theta_{i,\mathrm{off}} .
\label{eq:piper-dh-transform}
\end{equation}

The base-to-link-6 forward kinematics are

\begin{equation}
{}^{0}\!T_6(\mathbf q)=\prod_{i=1}^{6}{}^{i-1}\!A_i(q_i).
\label{eq:piper-fk}
\end{equation}

Table~\ref{tab:piper-dh} lists the parameters used by
\texttt{piper\_arm.py}. Lengths are in metres; angular values are shown as
multiples of $\pi$ or in degrees to match the implemented constants.

\begin{longtable}{@{}ccccc@{}}
\caption{Piper kinematic parameters implemented by the robot backend.}
\label{tab:piper-dh}\\
\toprule
$i$ & $\alpha_i$ & $a_i$ & $d_i$ & $\theta_{i,\mathrm{off}}$ \\
\midrule
\endfirsthead
\toprule
$i$ & $\alpha_i$ & $a_i$ & $d_i$ & $\theta_{i,\mathrm{off}}$ \\
\midrule
\endhead
\bottomrule
\endfoot
1 & $0$ & $0$ & $0.12300$ & $0$ \\
2 & $-\pi/2$ & $0$ & $0$ & $-172.2135102^\circ$ \\
3 & $0$ & $0.28503$ & $0$ & $-102.7827493^\circ$ \\
4 & $\pi/2$ & $-0.02198$ & $0.25075$ & $0$ \\
5 & $-\pi/2$ & $0$ & $0$ & $0$ \\
6 & $\pi/2$ & $0$ & $0$ & $0$ \\
\end{longtable}

The task-space point used for grasping is the midpoint between the finger
pads, located $0.06$~m along the link-6 $z$ axis.  Candidate IK solutions are
accepted only when their forward-kinematics residuals satisfy

\begin{align}
e_p &= \left\|\mathbf p_{\mathrm{FK}}-\mathbf p_{\mathrm{target}}\right\|_2
      < 2\ \mathrm{mm}, \\
e_R &= \cos^{-1}\!\left(\frac{\operatorname{tr}
      (R_{\mathrm{FK}}^{T}R_{\mathrm{target}})-1}{2}\right)
      < 0.01\ \mathrm{rad}.
\label{eq:piper-ik-residuals}
\end{align}

The original analytic solver evaluated only one shoulder/elbow branch.  The
revised solver enumerates shoulder, elbow, and wrist branches, rejects joint
limit violations, verifies every candidate using Eq.~\ref{eq:piper-ik-residuals},
and applies a damped DH-Jacobian numerical refinement when no closed-form
candidate passes.  A 500-case reachable-pose regression improved from 79\%
for the original single branch, to 89\% after branch enumeration, and to
483/500 (96.6\%, reported as 97\%) with the numerical fallback.  Because this
solver depends on the DH model rather than MuJoCo, the same implementation can
be used by both the DH-matched simulation model and the future physical Piper
planner.

The \texttt{piper\_real} MuJoCo model was generated from the official Piper
URDF and mesh assets in the ROS~Noetic branch of \texttt{piper\_ros}
(Reference~[16]). Across the checked configurations, the DH and MuJoCo link-6
frames agreed to within 0.09~mm. This removes the approximately 10~mm
kinematic mismatch observed when the community Menagerie model was driven by
the real-robot DH parameters (Reference~[18]).

\subsubsection{Motion-plan contract and execution semantics}

The planner returns a dependency-light \texttt{MotionPlan} rather than
issuing actuator commands directly. Its public structure is

\begin{center}
\texttt{start\_q $\rightarrow$ segments $\rightarrow$ waypoints[].q}.
\end{center}

Each segment has an explicit phase, including \texttt{TRANSIT},
\texttt{PRE\_APPROACH}, \texttt{FINAL\_APPROACH}, \texttt{CLOSING}, and
\texttt{LIFT}.  The plan also stores collision provenance through
\texttt{CollisionTrace}, including the phase, segment, interpolation sample,
contacting geometry, collision category, and rejection reason. Collision
checking includes segment endpoints and adaptively subdivided samples, with a
maximum joint interpolation step of 0.05~rad.

Planning and execution use the same saved joint waypoints. The executor does
not silently re-solve IK or skip a waypoint after validation. It rejects a
plan before motion when the observed start state differs from
\texttt{motion\_plan.start\_q}, when a segment seam is discontinuous, or when
the plan contains an invalid joint vector. This preserves the meaning of the
collision and feasibility checks performed during planning.

At task level, \texttt{TaskExecutor} implements four bounded operations:
\texttt{pick}, \texttt{place\_at}, \texttt{place\_into}, and
\texttt{clear\_table}.  A task progresses through an explicit finite-state
machine. For example, a container placement uses

\begin{equation*}
\begin{split}
\texttt{DETECT\_SOURCE} &\rightarrow \texttt{PLAN\_PICK}
\rightarrow \texttt{EXECUTE\_PICK}
\rightarrow \texttt{VERIFY\_GRASP} \\
&\rightarrow \texttt{DETECT\_TARGET}
\rightarrow \texttt{PLAN\_PLACE}
\rightarrow \texttt{EXECUTE\_PLACE}
\rightarrow \texttt{VERIFY\_DONE}.
\end{split}
\end{equation*}

Every state returns a structured result containing success, message, reason,
and duration. A failed grasp can trigger only a bounded recovery loop; the
pipeline otherwise stops at the first failed state and preserves the failure
provenance in JSON and CSV evidence.

\subsection{MuJoCo and Visual-Grasp Robot Backend}
\label{sec:visual-grasp-backend}

\subsubsection{Perception and simulation architecture}

The simulation uses MuJoCo~3.10.0 (Reference~[8]) with a 2~ms physics timestep
(500~Hz).  The
scene contains the Piper arm, a table, and textured cup, bottle, and bowl
models based on the YCB object set (Reference~[19]). Two cameras are supported:
a fixed world camera for eye-to-hand
perception and a wrist-mounted camera for eye-in-hand perception.  Both expose
known extrinsic transformations in simulation; the physical system will
replace these known transforms with independently verified camera
calibration.

The backend provides two perception modes with deliberately different roles:

\begin{itemize}
\item \texttt{sim\_gt} reads the object pose from MuJoCo and isolates the
      manipulation, IK, planning, and execution chain from perception error;
\item \texttt{yolo} detects the object in rendered RGB, uses the aligned depth
      image to estimate a three-dimensional centre, and transforms the centre
      into the robot base frame. A YOLO failure is reported as a failure and
      never falls back silently to \texttt{sim\_gt}.
\end{itemize}

The command bridge validates and normalizes the external JSON before
dispatch. The principal fields are shown in Table~\ref{tab:backend-command}.

\begin{longtable}{@{}
>{\raggedright\arraybackslash}p{0.25\linewidth}
>{\raggedright\arraybackslash}p{0.66\linewidth}@{}}
\caption{Implemented \texttt{visual\_grasp.bridge.v1} command fields.}
\label{tab:backend-command}\\
\toprule
Field & Implemented meaning \\
\midrule
\endfirsthead
\toprule
Field & Implemented meaning \\
\midrule
\endhead
\bottomrule
\endfoot
\texttt{task} & \texttt{pick}, \texttt{place\_at},
\texttt{place\_into}, or \texttt{clear\_table} \\
\texttt{source\_label}, \texttt{container\_label},
\texttt{labels}, \texttt{target\_xyz} & Task-dependent object and target
arguments with required-field and shape validation \\
\texttt{backend} & \texttt{sim\_mujoco} or \texttt{real\_piper} \\
\texttt{model} & Menagerie model or the DH-matched
\texttt{piper\_real} model \\
\texttt{perception} & \texttt{yolo} or \texttt{sim\_gt};
\texttt{real\_piper} rejects \texttt{sim\_gt} \\
\texttt{policy}, \texttt{speed} & Bounded grasp-policy and motion-speed
selection \\
\texttt{armed} & First explicit hardware opt-in; physical dispatch also
requires an independent API/CLI authorization \\
\end{longtable}

Dry-run validation is intentionally dependency-light: it does not import
MuJoCo, load a camera, construct a Piper SDK object, connect CAN, or move the
robot. Consequently, a successful dry run proves only schema validation and
dispatch selection.

\subsubsection{Quantitative simulation evidence}

Table~\ref{tab:backend-fixed-benchmark} reports a fixed-case
\texttt{place\_into} benchmark for a cup at $(x,y)=(0.38,-0.11)$~m. Both
backends used the DH-matched model and random seed 20260712. Task duration is
measured on the host from the first task state to the terminal task state. It
is not OpenHarmony, HDC, DDS, or ROS-bridge latency.

\begin{longtable}{@{}
>{\raggedright\arraybackslash}p{0.29\linewidth}rr@{}}
\caption{Fixed-case 100-trial cup placement benchmark.}
\label{tab:backend-fixed-benchmark}\\
\toprule
Metric & \texttt{sim\_gt} & YOLO/RGB-D \\
\midrule
\endfirsthead
\toprule
Metric & \texttt{sim\_gt} & YOLO/RGB-D \\
\midrule
\endhead
\bottomrule
\endfoot
Successful trials & 100/100 & 100/100 \\
Mean task duration (ms) & 1300.09 & 1857.06 \\
P95 task duration (ms) & 1322.52 & 1895.53 \\
Mean detection duration (ms) & 0.20 & 559.45 \\
P95 detection duration (ms) & 0.30 & 576.47 \\
Mean source-centre error (mm) & 0.00 & 6.10 \\
Mean lift height (mm) & 127.50 & 127.40 \\
Mean placement error (mm) & 29.50 & 28.30 \\
\end{longtable}

Because a single fixed pose does not establish workspace robustness, the
backend was also evaluated on bottle-position grids. A successful grasp
required a lift of at least 30~mm. Table~\ref{tab:backend-grid-benchmark}
reports the two relevant suites separately so that their different domains
and denominators are not conflated.

\begin{longtable}{@{}
>{\raggedright\arraybackslash}p{0.32\linewidth}
>{\raggedright\arraybackslash}p{0.16\linewidth}
>{\raggedright\arraybackslash}p{0.17\linewidth}
>{\raggedright\arraybackslash}p{0.25\linewidth}@{}}
\caption{Bottle perception-to-grasp grid evidence.}
\label{tab:backend-grid-benchmark}\\
\toprule
Suite & \texttt{sim\_gt} & YOLO/RGB-D & Interpretation \\
\midrule
\endfirsthead
\toprule
Suite & \texttt{sim\_gt} & YOLO/RGB-D & Interpretation \\
\midrule
\endhead
\bottomrule
\endfoot
Calibrated perception-to-grasp grid, 200 YOLO trials &
250/250 (100\%) & 191/200 (95.5\%) &
The calibrated suite passed its stated target; the YOLO Wilson 95\% interval
was 91.67--97.61\%. All nine failures occurred in
\texttt{DETECT\_SOURCE}. \\
Expanded $5\times5$ grid, ten repeats per cell &
250/250 (100\%) & 184/250 (73.6\%) &
Exact localization isolated a reliable manipulation chain, while the broader
YOLO evaluation did not meet the 85\% target. Its Wilson 95\% interval was
67.81--78.68\%, with 66 failures at \texttt{DETECT\_SOURCE}. \\
\end{longtable}

These results support a narrower conclusion than ``the complete visual grasp
system is 100\% successful.'' Under exact simulated localization, the
DH-matched IK, grasp planning, and MuJoCo execution chain completed all 250
expanded-grid trials. The broader camera-driven result was limited by target
detection and operating-domain rejection rather than by a demonstrated IK or
trajectory failure. Therefore, perception robustness remains open work and
the 73.6\% expanded-grid result must not be hidden by the 100/100 fixed-case
demonstration.

The machine-readable evidence used for these claims is:

\begin{itemize}
\item \path{multitask/benchmark_results_presentation/real_sim_gt_100_summary.json}
      and its trial CSV;
\item \path{multitask/benchmark_results_presentation/real_yolo_100_summary.json}
      and its trial CSV;
\item \path{multitask/benchmark_results/perception_e2e_v1/perception_e2e_summary.json}
      and its trial CSV;
\item \path{multitask/benchmark_results/bottle_heatmap_final_v8/summary.json}
      and its trial CSV.
\end{itemize}

\subsubsection{Physical Piper execution boundary}

The physical-control seam is implemented by a dependency-injected real-task
executor and a feedback-guarded \texttt{PiperHardwareBackend}. It reuses the neutral
\texttt{MotionPlan} contract and executes continuous joint-waypoint
polylines without re-solving IK. The backend rejects missing or stale
feedback, joint-limit violations, non-finite commands, start-state mismatch,
segment discontinuity, excessive requested or observed speed, following
error, robot fault/stop state, and waypoint timeout. Gripper actuation is an
explicit command and is not hidden inside ordinary waypoint execution.

Physical dispatch uses two independent opt-in gates: the normalized command
must contain \texttt{armed=true}, and the calling API or CLI must separately
set \texttt{allow\_hardware=true}. These gates issue an opaque authorization
capability that is required at the backend and SDK-adapter boundaries. Import,
construction, and dry-run paths do not connect hardware.

In the current checkout, 24 bridge-contract tests and 43 Piper-backend safety
tests pass (67/67 total). These tests use fakes or SDK-shaped stubs and verify
the software contract, cleanup, authorization, feedback, trajectory, and
fail-closed behaviour. They do \emph{not} prove CAN communication, physical
camera calibration, collision clearance around a real table, or safe motion
of the laboratory Piper. Moreover, the default SDK path cannot prove fresh
timestamps independently for all three joint pairs and therefore fails closed
with \texttt{JOINT\_FRESHNESS\_UNPROVEN}. Physical-arm acceptance remains
pending until the SDK/firmware interface, independent feedback freshness,
low-speed motion, emergency stop, and calibrated perception have been tested
on the actual platform.

\subsubsection{Backend-specific completion statement}

For DR3, the defensible status is therefore:

\begin{quote}
The \texttt{visual\_grasp} robot-operation backend is implemented and
quantitatively validated in MuJoCo, including DH-matched IK, bounded task
execution, structured command validation, and machine-readable benchmark
evidence. The feedback-guarded physical Piper execution seam is implemented
and passes offline safety tests, but physical-arm operation has not been
accepted and is not claimed as complete. The latest broad YOLO grid remains
below the 85\% target, so camera-driven perception robustness is retained as
an explicit limitation rather than being reported as a completed result.
\end{quote}
